% Laget automatisk av elevark/lag_elevark.py fra sim/friksjon.html. Endringer her blir overskrevet.
\documentclass[11pt]{article}
\usepackage{elevark}
\begin{document}
\elevarkhode{Fysikk 1 · Krefter}{Friksjon}{https://fysikklab.no/sim/friksjon.html}{Simuleringen}

Når du drar i en kloss som ligger på et bord, står den først i ro. Friksjonen holder igjen like hardt som du drar, opp til en grense. Så glipper klossen, og da trengs det mindre kraft for å holde den i gang. Hvordan henger dette sammen med normalkraften?

\section*{Oppgaver}

Start med tre mot tre, 1,0~kg og økende kraft.

\begin{oppgaveliste}
\oppgave \fase{gjett}{Gjett}Hvordan tror du grafen for friksjonskraften ser ut mens du drar sakte og stadig hardere? Kjør simuleringen og sammenlign.
\svar
\oppgave \fase{forklar}{Forklar}Mens klossen står i ro, er friksjonskraften like stor som trekkraften. Hvorfor? Bruk Newtons 1. lov.
\svar
\oppgave \fase{test}{Regn}Regn ut normalkraften, største hvilefriksjon og glidefriksjonen for klossen. Ved hvilken trekkraft glipper den? Sjekk.
\svar
\oppgave \fase{test}{Undersøk}Doble massen. Hva skjer med kraften som skal til for at klossen skal glippe? Hva skjer med forholdet mellom største hvilefriksjon og glidefriksjon?
\svar
\oppgave \fase{test}{Regn}Velg konstant kraft. Klossen på 1,0~kg dras med 4,0~N og så med 6,0~N. Hva skjer i hvert tilfelle? Regn ut akselerasjonen når klossen glir.
\svar
\oppgave \fase{forklar}{Vurder}I forsøket med hvilefriksjon og glidefriksjon målte dere forholdet mellom dem for ulike flater. Hva er forholdet for flatene i simuleringen? Er det det samme for alle?
\svar
\oppgave \fase{utfordring}{Utfordring}En bil bremser med låste hjul på tørr asfalt, og en annen bremser med ABS, der hjulene ruller helt på grensen til å skli. Hvilken stopper raskest? Bruk forskjellen på hvilefriksjon og glidefriksjon.
\svar
\end{oppgaveliste}

\end{document}
